\begin{figure}[htbp]
\centering
\begin{tikzpicture}[
  node distance = 0.5cm,
  font = \small,
  >=Stealth
]

\tikzset{
  block/.style  = {rectangle, draw, rounded corners=3pt, fill=#1,
                   minimum width=2.8cm, minimum height=1.0cm,
                   text width=2.6cm, align=center},
  sumnode/.style = {circle, draw, fill=white, minimum size=0.65cm,
                    inner sep=0pt, font=\normalsize},
  clipb/.style  = {rectangle, draw, fill=gray!12, rounded corners=3pt,
                   minimum width=1.8cm, minimum height=0.8cm,
                   text width=1.6cm, align=center, font=\footnotesize},
  signal/.style = {font=\footnotesize},
  arr/.style    = {->, thick},
  fdbk/.style   = {->, thick, draw=gray!70, dashed}
}

%% ─── INPUT LABELS ───────────────────────────────────────────
\node[signal] (sin_mamba) at (-4.2,  1.5) {$\mathbf{s}_{t-W:t}\;\in\mathbb{R}^{18\times W}$};
\node[signal] (sin_ppo)   at (-4.2, -1.5) {$\mathbf{s}_t\;\in\mathbb{R}^{18}$};
\node[signal] (sin_eta)   at (-4.2, -2.5) {$\hat{\boldsymbol{\eta}}\;\in\mathbb{R}^{4}$};

%% ─── CONTROLLER BLOCKS ──────────────────────────────────────
\node[block=blue!12] (mamba) at (0, 1.5)
    {Mamba SSM\\{\scriptsize controlador nominal\\2L · $d_{\rm model}$=128}};

\node[block=orange!15] (ppo) at (0, -2.0)
    {PPO\\{\scriptsize compensador\\de fallas}};

%% ─── SUM NODE ───────────────────────────────────────────────
\node[sumnode] (suma) at (3.8, 0) {$+$};

%% ─── CLIP BOX ───────────────────────────────────────────────
\node[clipb] (clip) at (6.0, 0)
    {clip\\$[\omega_{\min},\omega_{\max}]$};

%% ─── MOTORS / FAULT ─────────────────────────────────────────
\node[block=red!8] (motors) at (8.5, 0)
    {Motores\\{\scriptsize $\omega_i^{\rm eff}=\eta_i\cdot u_{t,i}$}};

%% ─── UAV ────────────────────────────────────────────────────
\node[block=blue!18] (uav) at (11.5, 0)
    {UAV\\{\scriptsize CF2X\\gym-pybullet}};

%% ─── OUTPUT / LABEL NODES ───────────────────────────────────
\node[signal] at (2.15, 2.05)  {$\mathbf{u}_t^{\rm nom}\!\in\!\mathbb{R}^4$};
\node[signal] at (2.15, -1.55) {$\Delta\mathbf{u}_t\!\in\!\mathbb{R}^4$};
\node[signal] at (4.85, 0.28)  {$\mathbf{u}_t$};

%% ─── ARROWS: inputs → blocks ────────────────────────────────
\draw[arr] (sin_mamba.east) -- (mamba.west);
\draw[arr] (sin_ppo.east)   -- (ppo.west);
\draw[arr] (sin_eta.east)   -- ++(0.5,0) |- (ppo.west);

%% ─── ARROWS: blocks → sum ───────────────────────────────────
\draw[arr] (mamba.east) -- ++(0.6,0) |- (suma.north);
\draw[arr] (ppo.east)   -- ++(0.6,0) |- (suma.south);

%% ─── ARROWS: sum → clip → motors → UAV ─────────────────────
\draw[arr] (suma)   -- (clip);
\draw[arr] (clip)   -- (motors);
\draw[arr] (motors) -- (uav);

%% ─── FEEDBACK: UAV → inputs ─────────────────────────────────
%% Bring s_t out from UAV bottom, route below everything, back to inputs
\draw[fdbk]
    (uav.south) -- ++(0,-1.2)
    -- ++(-15.7,0)
    -- ++(0, 1.0)
    -- (sin_mamba.west);

\draw[fdbk]
    (uav.south) ++(0,-1.2) -- ++(-15.7,0) -- ++(0,0.8)
    -- (sin_ppo.west);

%% ─── FAULT ANNOTATION ───────────────────────────────────────
\node[signal, text=red!70, above=0.15cm of motors]
    {\scriptsize falla $\eta_i\!\in\![0,1]$};

%% ─── BRACE / BOX around the two-level system ────────────────
\begin{scope}[on background layer]
  \node[draw=blue!30, dashed, rounded corners=6pt,
        fill=blue!3, inner sep=8pt,
        fit=(mamba)(ppo)(suma)(clip)] (box) {};
\end{scope}
\node[above=2pt of box, font=\scriptsize\itshape, text=blue!60]
    {Sistema FTC propuesto};

\end{tikzpicture}
\caption{Arquitectura de control tolerante a fallas de dos niveles propuesta en esta investigación. El controlador Mamba SSM genera la acción nominal $\mathbf{u}_t^{\rm nom}$ mediante aprendizaje por imitación del PID; el compensador PPO genera la corrección $\Delta\mathbf{u}_t$ condicionada al estado del sistema y al vector de efectividad estimado $\hat{\boldsymbol{\eta}}$. La señal resultante $\mathbf{u}_t$ se recorta al rango físico de los motores antes de ejecutarse. La trayectoria retroalimentada $\mathbf{s}_t$ proviene del estado del UAV en cada paso de control.}
\label{fig:arquitectura}
\end{figure}
